\chapter{Background}
\label{chp:background}

Anonymous communication aims to enable users to exchange messages over a network without revealing who is communicating with whom, thereby protecting both sender and receiver identities. Many such systems have been proposed and implemented over the years, such as Tor, I2P, and mixnets. Most of these designs depend on the \emph{threshold model} to provide anonymity, requiring at least some components of their infrastructure to behave honestly. In practice, this assumption is increasingly strained by large-scale surveillance, correlation attacks, and the growing centralization of internet infrastructure~\cite{sun:2015:RAPTOR}.

Systems with stronger anonymity guarantees based on \gls{mpc} exist. However, these suffer from poor scalability due to their high interactivity, making them impractical for large-scale networks. A \gls{dc-net}, for example, can provide anonymity even if there are no trusted components in the network infrastructure. Doing so requires every participant to interact with every other participant, for every round of messages sent.

Recently, a new line of research has proposed fundamentally different designs, aiming to offer strong anonymity even in the presence of compromised or adversarial components in the network actively attempting to break anonymity. The first such design was the \gls{niar}~\cite{cryptoeprint:2021/435}, which proposing using \gls{mcfe} to allow a single, untrusted router to shuffle or permute messages from a set of clients \emph{obliviously}; without learning anything about the permutation. Thereby making it infeasible for the router to link any of the messages to its sender.

\gls{niar} was a breakthrough in anonymous communication, not only showing the theoretical possibility of designing anonymous routers, but also that such routers can be \emph{non-interactive}: clients only send one message to the router and do not need to perform any additional interaction to achieve anonymity. Using a centralized (anonymous) router in this way yields similar benefits to existing \gls{mpc}-based systems, without needing any interactions between clients. However, there are a number of problems that makes \gls{niar} impractical, if not impossible, to implement in practice.

\begin{enumerate}
    \item The cryptographic components required by NIAR are extremely computationally expensive or not known to be implementable in practice.
    \item Achieving anonymity for recipients assumes the existence of a very strong primitive called \emph{indistinguishability obfuscation with sub-exponential security}: guaranteeing that no efficient adversary can distinguish between the obfuscations of two functionally equivalent circuits with more than a sub-exponentially small advantage. The practicality of such a construction is unknown.
    \item Reusing the same setup across multiple rounds makes messages from the same sender linkable, unless the expensive setup phase is rerun for every round.
\end{enumerate}

These problems were highlighted in \gls{spar}

\begin{itemize}
    \item Anonymous communication
    \item Current solutions
    \item Problems with current solutions
    \item \gls{fhe} as a new approach
    \item FHE 'holy grail' of privacy-preserving communication
    \item \gls{spar}
    \item RSA is multiplicatively homomorphic, but (follow-up) paper theorized a \emph{fully} homomorphic scheme may be possible \cite{gentry:2009, rivest:1978:databanks}
    \item \emph{Contributions}
    \begin{itemize}
        \item Port of TFHE-style external/internal products to BGV-RNS
        \item A reusable OpenFHE library
        \item Complete sPAR implementation
        \item First empirical evaluation of it
    \end{itemize}
\end{itemize}

\subsubsection{Structure of the thesis}
\begin{itemize}
    \item Building blocks and related works; FHE, gadget decomposition, TFHE, etc.
    \item Explanation of the code
    \item Experimental setup i.e. how the numbers where found
    \item Results, benchmarks, noise analysis 
    \item Future work
\end{itemize}

%% TODO: nowhere in the thesis is there an explicit "this thesis contributes:" list. You have at least four defensible contributions scattered across chapter intros: (a) the first (to your knowledge — verify) port of TFHE-style external/internal products to BGV-RNS, as a reusable OpenFHE library; (b) a complete sPAR implementation and the first empirical evaluation of it; (c) the noise-asymmetry ordering analysis and the resulting 30–40 bit improvement in Server::Write; (d) the hybrid-RGSW construction and its open problem. Claim them in the introduction, in this order of confidence, and map each to the chapter that delivers it.